\documentclass[10pt]{beamer}

% Tema Limpo e Institucional IFCE Taua
\usetheme{Madrid}
\usecolortheme{whale}

% Pacotes Essenciais
\usepackage[utf8]{inputenc}
\usepackage[T1]{fontenc}
\usepackage[brazilian]{babel}
\usepackage{graphicx}
\usepackage{booktabs}
\usepackage{amsmath}
\usepackage{amssymb}
\usepackage{listings}
\usepackage{xcolor}
\usepackage{tikz}
\usetikzlibrary{shapes.geometric,arrows.meta,positioning,calc}

% Configuracao de Cores e Codigo (Python)
\definecolor{codegreen}{rgb}{0,0.6,0}
\definecolor{codegray}{rgb}{0.5,0.5,0.5}
\definecolor{codepurple}{rgb}{0.58,0,0.82}
\definecolor{backcolour}{rgb}{0.96,0.97,0.98}

\lstset{
    language=Python,
    backgroundcolor=\color{backcolour},
    basicstyle=\ttfamily\scriptsize,
    keywordstyle=\color{blue}\bfseries,
    stringstyle=\color{red!80!black},
    commentstyle=\color{green!60!black},
    numberstyle=\tiny\color{codegray},
    breaklines=true,
    showstringspaces=false,
    tabsize=4,
    frame=single,
    rulecolor=\color{blue!20}
}

% Metadados da Apresentacao
\title[Engenharia de Software -- Aula 27]{Engenharia de Software -- Aula 27}
\subtitle{Sprint 04: Testes Automatizados (Pytest), CLI e Fechamento da Sprint}
\author[Prof. Me. Reginaldo Fernandes]{Prof. Me. Reginaldo Pereira Fernandes}
\institute[IFCE Campus Tauá]{
    Instituto Federal de Educação, Ciência e Tecnologia do Ceará -- Campus Tauá\\
    Técnico em Informática para Internet (TSII.210)
}
\date{Semana 2 -- Encontro 4 (100 min)}

\begin{document}

% Slide 1: Titulo
\begin{frame}
    \titlepage
\end{frame}

% Slide 2: Agenda do Encontro
\begin{frame}{Agenda do Encontro 4 (100 Minutos)}
    \tableofcontents
\end{frame}

\section{Garantia de Qualidade e Pirâmide de Testes}

\begin{frame}{Por que Automatizar Testes de Software?}
    \begin{itemize}
        \item \textbf{A Ilusão do Teste Manual:}
        \begin{itemize}
            \item Testar "clicando na tela" é lento, não repetível e propenso a esquecimentos.
            \item Toda alteração no código exige re-testar tudo do zero (regressão!).
        \end{itemize}
        \item \textbf{A Filosofia do Pytest:}
        \begin{itemize}
            \item Testes são código que valida código.
            \item Execução de centenas de cenários em poucos milissegundos.
            \item Confiança total antes de enviar o Pull Request para a liderança.
        \end{itemize}
    \end{itemize}
\end{frame}

\begin{frame}{Diagrama da Pirâmide de Testes e Rastreabilidade}
\centering
\resizebox{0.92\textwidth}{!}{%
\begin{tikzpicture}[>=Latex, font=\sffamily]
    % Piramide
    \draw[thick, draw=red!70!black, fill=red!15] (5, 3.2) -- (3.8, 1.8) -- (6.2, 1.8) -- cycle;
    \node[font=\tiny\bfseries] at (5, 2.3) {E2E / CLI};

    \draw[thick, draw=yellow!80!black, fill=yellow!20] (3.8, 1.8) -- (2.6, 0.4) -- (7.4, 0.4) -- (6.2, 1.8) -- cycle;
    \node[font=\tiny\bfseries] at (5, 1.1) {Integração (SQLite em Memória)};

    \draw[thick, draw=green!70!black, fill=green!20] (2.6, 0.4) -- (1.0, -1.2) -- (9.0, -1.2) -- (7.4, 0.4) -- cycle;
    \node[font=\scriptsize\bfseries] at (5, -0.4) {Testes Unitários (DUC: Regras \& Exceções)};

    % Rastreabilidade DUC
    \node[anchor=west, font=\tiny, text=blue!80!black] at (9.2, -0.4) {$\leftarrow$ Valida FE01..FE04, FA01, FA02};
    \node[anchor=west, font=\tiny, text=blue!80!black] at (7.6, 1.1) {$\leftarrow$ Valida Baixa no Estoque e Banco};
    \node[anchor=west, font=\tiny, text=blue!80!black] at (6.4, 2.3) {$\leftarrow$ Valida Interface de Terminal};
\end{tikzpicture}
}
\end{frame}

\section{Git Flow e Ciclo de Vida da Sprint}

\begin{frame}{Diagrama de Git Flow: Da Feature ao Pull Request}
\centering
\resizebox{0.95\textwidth}{!}{%
\begin{tikzpicture}[>=Latex, font=\sffamily]
    % Main branch
    \draw[very thick, draw=blue!80!black] (0, 1.8) -- (10, 1.8);
    \node[anchor=east, font=\scriptsize\bfseries, text=blue!80!black] at (-0.2, 1.8) {main};
    \node[draw, circle, fill=blue!80!black, inner sep=1.5pt] at (1, 1.8) {};
    \node[draw, circle, fill=blue!80!black, inner sep=1.5pt] at (9, 1.8) {};
    \node[above=0.1cm, font=\tiny] at (1, 1.8) {v0.3.0};
    \node[above=0.1cm, font=\tiny\bfseries] at (9, 1.8) {v0.4.0 (Merge)};

    % Feature branch
    \draw[very thick, draw=green!60!black] (1, 1.8) -- (2, 0) -- (8, 0) -- (9, 1.8);
    \node[anchor=east, font=\scriptsize\bfseries, text=green!60!black] at (1.8, -0.4) {feature/task-204};

    % Commits semanticos
    \node[draw, circle, fill=green!60!black, inner sep=1.5pt] (c1) at (3, 0) {};
    \node[draw, circle, fill=green!60!black, inner sep=1.5pt] (c2) at (4.5, 0) {};
    \node[draw, circle, fill=green!60!black, inner sep=1.5pt] (c3) at (6, 0) {};
    \node[draw, circle, fill=green!60!black, inner sep=1.5pt] (c4) at (7.5, 0) {};

    \node[below=0.1cm of c1, font=\tiny, align=center] {feat(domain):\\modelo};
    \node[below=0.1cm of c2, font=\tiny, align=center] {feat(db):\\sqlite repo};
    \node[below=0.1cm of c3, font=\tiny, align=center] {feat(service):\\checkout};
    \node[below=0.1cm of c4, font=\tiny, align=center] {test(pytest):\\10 testes};
\end{tikzpicture}
}
\vspace{0.3cm}
\begin{alertblock}{Conventional Commits na Equipe}
Commits atômicos e claros facilitam o code review e a rastreabilidade do projeto!
\end{alertblock}
\end{frame}

\section{Implementação: Suíte de Testes com Pytest}

\begin{frame}[fragile]{Mão na Massa: Fixture com SQLite em Memória (\texttt{test\_cupons.py})}
\begin{lstlisting}[language=Python]
import pytest
from datetime import date, timedelta
from bibliotaua.database import obter_conexao, SCHEMA_SQL, popular_dados_iniciais
from bibliotaua.repositories.cupom_repository import CupomRepository
from bibliotaua.models.cupom import Cupom, TipoDesconto

@pytest.fixture
def conn():
    """Cria um banco SQLite limpo em memoria para cada teste."""
    conexao = obter_conexao(":memory:")
    conexao.executescript(SCHEMA_SQL)
    popular_dados_iniciais(conexao)
    yield conexao
    conexao.close()

@pytest.fixture
def cupom_repo(conn):
    return CupomRepository(conn)
\end{lstlisting}
\end{frame}

\begin{frame}[fragile]{Mão na Massa: Testando Caminho Feliz e Exceções}
\begin{lstlisting}[language=Python]
def test_aplicar_cupom_percentual_sucesso(cupom_repo):
    cupom = cupom_repo.buscar_por_codigo("VOLTAAULAS10")
    # Livro de R$ 100.00 -> 10% de desconto = R$ 10.00
    desconto = cupom.calcular_desconto(100.0)
    assert desconto == 10.0

def test_cupom_expirado_dispara_excecao(conn, cupom_repo):
    cupom = Cupom(
        codigo="EXPIRADO",
        tipo_desconto=TipoDesconto.VALOR_FIXO,
        valor_desconto=15.0,
        validade_ate=date.today() - timedelta(days=1)
    )
    cupom_repo.salvar(cupom)
    with pytest.raises(CupomExpiradoError):
        cupom.calcular_desconto(50.0)
\end{lstlisting}
\end{frame}

\begin{frame}[fragile]{Execução do Pytest: 100\% Verde no Terminal}
\begin{lstlisting}[language=bash]
$ PYTHONPATH=. pytest bibliotaua/tests/test_cupons.py -v

test_cupons.py::test_buscar_cupom_existente PASSED              [ 10%]
test_cupons.py::test_buscar_cupom_inexistente_retorna_none PASSED [ 20%]
test_cupons.py::test_aplicar_cupom_percentual_sucesso PASSED    [ 30%]
test_cupons.py::test_aplicar_cupom_valor_fixo_sucesso PASSED    [ 40%]
test_cupons.py::test_cupom_inativo_dispara_excecao PASSED       [ 50%]
test_cupons.py::test_cupom_expirado_dispara_excecao PASSED      [ 60%]
test_cupons.py::test_cupom_esgotado_dispara_excecao PASSED      [ 70%]
test_cupons.py::test_subtotal_abaixo_do_minimo_dispara_excecao PASSED [ 80%]
test_cupons.py::test_checkout_completo_com_cupom PASSED         [ 90%]
test_cupons.py::test_decremento_estoque_livros_no_checkout PASSED [100%]

============================== 10 passed in 0.03s ==============================
\end{lstlisting}
\end{frame}

\section{Demonstração Interativa (CLI)}

\begin{frame}[fragile]{Homologação no Terminal: \texttt{bibliotaua/main.py}}
\begin{lstlisting}[language=bash]
$ python3 -m bibliotaua.main
==================================================
   BIBLIOTAUA - LIVRARIA DO IFCE CAMPUS TAUA
==================================================
1. Acervo de Livros Disponiveis
2. Simular Compra no Carrinho
3. Sair
Escolha uma opcao: 2

Livro: Introducao ao Python (Qtd: 1) - R$ 80.00
Cupom Informado: VOLTAAULAS10
--------------------------------------------------
Subtotal:          R$ 80.00
Desconto (10%):  - R$  8.00
Total a Pagar:     R$ 72.00 (PIX Gerado com Sucesso!)
--------------------------------------------------
Estoque atualizado de 15 para 14 exemplares.
\end{lstlisting}
\end{frame}

\section{Fechamento da Sprint e Pull Request}

\begin{frame}{Rito Ágil: Submissão do Pull Request (PR)}
    \begin{block}{Template de PR para a Tech Lead Aline Mendes}
    \begin{itemize}
        \item \textbf{Título:} \texttt{feat(cupons): Motor de Validação e Aplicação (TASK-204)}
        \item \textbf{Descrição do que foi feito:}
        \begin{itemize}
            \item Mapeamento das regras da DUC UC04 e classes UML.
            \item Tabela \texttt{cupons} criada no SQLite com constraints estritas.
            \item \texttt{CupomRepository} implementado com proteção contra SQL Injection.
            \item \texttt{CheckoutService} coordenando faturamento e baixa de estoque.
            \item 10 testes no Pytest com 100\% de cobertura dos fluxos de exceção.
        \end{itemize}
        \item \textbf{Status da Entrega:} Homologado e Pronto para Merge!
    \end{itemize}
    \end{block}
\end{frame}

\begin{frame}{Síntese da Sprint 04: O Que Construímos em 4 Encontros?}
    \begin{itemize}
        \item \textbf{Encontro 1:} Diagrama de Casos de Uso, Classes, DUC, Exceções de Domínio e Entidade \texttt{Cupom}.
        \item \textbf{Encontro 2:} Diagrama Relacional (DER), script DDL SQLite e \texttt{CupomRepository}.
        \item \textbf{Encontro 3:} Diagrama de Sequência, \texttt{CalculadoraDescontoService} e \texttt{CheckoutService}.
        \item \textbf{Encontro 4:} Testes automatizados com Pytest, demonstração CLI e Pull Request.
    \end{itemize}
    \vspace{0.4cm}
    \centering
    \textbf{Parabéns pelo trabalho profissional de Engenharia de Software!}
\end{frame}

\end{document}
